\documentclass[10pt,a4paper]{amsart}
\usepackage[margin=19mm]{geometry}
\usepackage{amsmath,amssymb,mathtools}
\usepackage[scr=rsfs,cal=euler]{mathalfa}
\usepackage{xfrac}
\usepackage[unicode,hidelinks,bookmarks=false]{hyperref}
\newcommand{\ie}{i.e.\ }
\newcommand{\cf}{cf.\ }
\newcommand{\resp}{respectively\ }
\newcommand{\Subsection}{\subsection}
\providecommand{\nobreakdash}{\nobreak\hbox{-}}
\setlength{\emergencystretch}{2em}
\multlinegap0pt
\usepackage{amssymb}
\usepackage{xcolor}
\usepackage{float}
\newtheorem{proposition}{Proposition}
\title{Invariant theory for non-reductive actions: extensions of Hilbert and Schwarz theorems}
\author{Leandro Nery}
\date{Source: arXiv:2510.19053v2 (CC BY 4.0)}
\begin{document}
\maketitle

\noindent\emph{Source and licence.} Invariant theory for non-reductive actions: extensions of Hilbert and Schwarz theorems. Version arXiv:2510.19053v2 (CC BY 4.0), distributed by arXiv under the Creative Commons Attribution 4.0 International licence, http://creativecommons.org/licenses/by/4.0/, as recorded in the arXiv licence metadata for this exact version and shown on the abstract page https://arxiv.org/abs/2510.19053v2. Full licence notice: \texttt{licenses/arxiv-2510.19053-CC-BY-4.0.txt}.\par\smallskip

\noindent\textit{Editorial scope.} Counted unit: the article's proposition prop:insufficiency (source0.tex lines 390--395). The article's complete original proof of it is delivered with it (source0.tex lines 397--409, one \texttt{proof} environment of 13 lines). No mathematics is written, repaired or re-derived: the statement and the proof are the article's own lines, in the article's own notation, and no retained line is substituted or re-flowed.  No further source-local context is needed: every cross-reference of every retained line resolves inside the two delivered spans.  Nothing was omitted from any retained span, so the package carries no omission record.  No retained line contains a citation, so the wrapper carries no bibliography and no bibliography entry.  No retained line contains a figure environment, an image-inclusion command or drawing code, so no image is delivered and the package declares no asset dependency.  Editorial formatting only, in addition: an AMS \texttt{amsart} preamble that restates the article's own declarations and macros used by the retained lines, namely the environments \texttt{proposition} on separate counters, together with the standard packages those lines need.  The retained environments print in this document as Proposition 1 in order of appearance, whereas the article prints the same items with the numbering of its own full document.  The complete native source of the article as distributed in the arXiv e-print for this exact version is bundled unchanged as \texttt{sources/187547/source0.tex}.
\begin{proposition} \label{prop:insufficiency}
	Let $\Gamma = \langle \mathcal{H}_\beta \rangle$ be the discrete subgroup of $O(1,1)$ generated by a hyperbolic rotation. Let $\rho(x,y) = x^2 - y^2$ be the generator of the invariant ring $\mathcal{P}(\mathbb{R}^{1,1})^\Gamma$. Then the pullback map $\rho^* : \mathcal{E}(\mathbb{R}) \to \mathcal{E}(\mathbb{R}^2)^\Gamma$ is not surjective; that is,
	\[
	\rho^*\mathcal{E}(\mathbb{R}) \subsetneq \mathcal{E}(\mathbb{R}^2)^\Gamma.
	\]
\end{proposition}

\begin{proof}
	We construct a smooth $\Gamma$-invariant function $G$ that is not a function of the quadratic form $\rho$. Let $G: \mathbb{R}^2 \to \mathbb{R}$ be defined by
	\[
	G(x,y) =
	\begin{cases}
		\exp\left(-\frac{1}{x^2 - y^2}\right), & \text{if } x^2 - y^2 > 0 \text{ and } x > 0, \\0, & \text{otherwise}.\end{cases}
	\]
	To verify that $G \in \mathcal{E}(\mathbb{R}^2)$, note that $G$ is clearly smooth on the open regions $\{x^2 - y^2 > 0, x > 0\}$ and its complement. Along the light-cone $x^2 - y^2 = 0$, the function and all its partial derivatives vanish as $t = \rho(x,y) \to 0^+$ due to the rapid decay of the exponential term, ensuring $C^\infty$ regularity across the boundary.
	
	The $\Gamma$-invariance of $G$ is best seen in light-cone coordinates $u = x+y$ and $v = x-y$, where the action is $(u,v) \mapsto (ue^\beta, ve^{-\beta})$. The region $\rho(x,y) > 0$ comprises two disjoint components: the right branch ($u, v > 0$) and the left branch ($u, v < 0$). Since the action of $\mathcal{H}_\beta$ preserves the quadrants, each branch is an invariant manifold. On the right branch, $G(u,v) = \exp(-1/uv)$, which is invariant under the scaling $u \mapsto ue^\beta$ and $v \mapsto ve^{-\beta}$. Since $G$ vanishes identically on the left branch and in the region $uv \le 0$, it follows that $G \in \mathcal{E}(\mathbb{R}^2)^\Gamma$.
	
	Finally, suppose there exists $f \in \mathcal{E}(\mathbb{R})$ such that $G = f \circ \rho$. This would imply that $G$ is constant on the level sets of $\rho$. However, for any fixed $t > 0$, the level set $\rho^{-1}(t)$ contains points in both the right and left branches. On the right branch, $G$ takes the value $e^{-1/t} > 0$, while on the left branch, $G$ is identically zero. This contradiction shows that $G \notin \rho^*\mathcal{E}(\mathbb{R})$, completing the proof.
\end{proof}

\end{document}
